\section{Synthesis for This Thesis}
\label{sec:background_synthesis}

The theoretical background can be summarized as a mapping problem. Virtual embodiment depends on how the user's physical body, actions, and sensory feedback are related to a virtual body or action. Agency, ownership, and self-location provide separate but connected ways of describing this relation. Visuomotor congruence, spatial alignment, temporal synchrony, and plausible body mapping are therefore central design variables.

In XR rowing, these variables are unusually concrete. The user performs a real rowing-like movement on a physical ergometer. The system must decide whether the virtual representation should remain close to this physical setup or transform it into a more complete on-water rowing scene. Both strategies are defensible, but they support different forms of realism.

The ergometer-congruent mode prioritizes correspondence with the user's real action. It is expected to support agency because the visible movement can follow the user's own stroke closely. It may also support ownership because the virtual hands, body, handle, and seat can be arranged in a spatially plausible relation to the user's felt body and apparatus. Its limitation is that it may feel less like the culturally recognizable image of rowing on water.

The boat-centred mode prioritizes sport depiction and environmental plausibility. It can show a rower in a shell, moving through water, using oars, and performing a visually recognizable rowing action. This may strengthen presence, self-location in the rowing scene, and the subjective meaning of the activity as rowing. Its limitation is that it may require a transformation between the felt ergometer handle path and the visible oar-based movement.

This thesis therefore does not assume that one representation is universally better. Instead, it expects that the two modes may produce different embodiment profiles. The ergometer-congruent mode may be stronger for direct agency and body correspondence. The boat-centred mode may be stronger for plausibility, sport realism, and the feeling of being situated in a rowing environment. Body ownership and self-location may depend on how participants integrate these cues.

The empirical comparison is built around this tension. By holding the general VR rowing context constant while varying the representation strategy, the study examines whether embodiment in XR rowing is more strongly supported by precise movement correspondence or by a more realistic sport depiction. The result should help clarify how embodiment concepts can inform the design of future XR exercise and sport systems.
